\documentclass[12pt,a4paper]{article}
\usepackage[utf8]{inputenc}
\usepackage[spanish]{babel}
\usepackage{cite}
\usepackage{longtable} 

\title {Taller de tipos de requisitos (FORM-TL-01) \\ 
\vspace{0.3cm} 
\large {Plataforma de gestión de citas médicas}}
\author{Cristopher Lavayen}
\date{7 de Noviembre de 2026}

\begin{document}
\maketitle

\section{Clasificación de Enunciados}

\renewcommand{\arraystretch}{1.5}
\begin{longtable}{|c|p{3.5cm}|c|c|p{5cm}|}
\caption{Clasificación de los enunciados E01 al E20}
\label{tab:clasificacion} \\
\hline
\textbf{Código} & \textbf{Tipo (ISO 25010)} & \textbf{Prod/Proc} & \textbf{Sist/Soft} & \textbf{Justificación} \\
\hline
\endfirsthead

\multicolumn{5}{c}{{\tablename\ \thetable{} -- Continuación de la página anterior}} \\
\hline
\textbf{Código} & \textbf{Tipo (ISO 25010)} & \textbf{Prod/Proc} & \textbf{Sist/Soft} & \textbf{Justificación} \\
\hline
\endhead

\hline
\endfoot

\hline
\endlastfoot


E01 & Funcional & Producto & Software & Describe una acción directa que el paciente ejecutará en la plataforma. \\ \hline
E02 & Funcional & Producto & Software & Define un comportamiento automático de notificación del aplicativo. \\ \hline
E03 & Calidad (Eficiencia de desempeño) & Producto & Software & Especifica una característica de rendimiento, redactada de forma vaga. \\ \hline
E04 & Funcional & Producto & Software & Detalla una interacción de ingreso de datos del recepcionista. \\ \hline
E05 & Restricción & Producto & Software & Impone una limitación tecnológica obligatoria sobre el stack de desarrollo. \\ \hline

E06 & Funcional & Producto & Software & Permite al médico visualizar información y actualizar el estado de una cita. \\ \hline
E07 & Calidad (Seguridad) & Producto & Software & Establece una protección de información, redactada vagamente. \\ \hline
E08 & Restricción & Producto & Software & Condiciona la compatibilidad tecnológica obligatoria de la app móvil. \\ \hline
E09 & Funcional & Producto & Software & Implementa una regla de negocio para penalizar cancelaciones tardías. \\ \hline
E10 & Calidad (Usabilidad) & Producto & Software & Define la facilidad de uso, redactada vagamente para adultos mayores. \\ \hline

E11 & Restricción & Proceso & Sistema & Establece un límite de tiempo (6 meses) para la ejecución del proyecto. \\ \hline
E12 & Funcional & Producto & Software & Define un comportamiento automático de cierre de sesión por inactividad. \\ \hline
E13 & Funcional & Producto & Sistema & Describe una interacción que involucra hardware físico (quiosco, impresora). \\ \hline
E14 & Calidad (Fiabilidad) & Producto & Software & Especifica una exigencia de disponibilidad del servicio, redactada vagamente. \\ \hline
E15 & Funcional & Producto & Software & Permite al director médico generar un reporte consolidado de inasistencias. \\ \hline

E16 & Restricción & Proceso & Sistema & Define el ciclo de vida iterativo y de validación del proyecto. \\ \hline
E17 & Restricción & Producto & Sistema & Impone cumplimiento normativo y legal sobre la plataforma. \\ \hline
E18 & Calidad (Compatibilidad) & Producto & Software & Exige interoperabilidad, redactada de forma vaga ("sin problemas"). \\ \hline
E19 & Calidad (Fiabilidad) & Producto & Software & Métrica específica de recuperación ante fallos del servidor. \\ \hline
E20 & Restricción & Proceso & Software & Establece una regla de calidad obligatoria en el flujo de desarrollo. \\
\end{longtable}

\section{Requisitos Cuantificados}

\section{Discrepancias y Resolución}

\bibliographystyle{ieeetr}
\bibliography{referencias}
\end{document}